\documentclass[aspectratio=169,10pt]{beamer}
\usetheme{default}
\usecolortheme{default}

\usepackage{graphicx}
\usepackage{adjustbox}
\usepackage{amsmath}
\usepackage{amssymb}
\usepackage{booktabs}
\usepackage{appendixnumberbeamer}

\setbeamertemplate{navigation symbols}{}
\setbeamertemplate{items}[circle]

\setbeamertemplate{footline}{
  \hbox{%
    \begin{beamercolorbox}[wd=.5\paperwidth,ht=2.5ex,dp=1ex,left]{author in head/foot}%
      \hspace*{2ex}\insertshortauthor, \insertshortinstitute\ (\insertdate)
    \end{beamercolorbox}%
    \begin{beamercolorbox}[wd=.5\paperwidth,ht=2.5ex,dp=1ex,right]{date in head/foot}%
      \insertframenumber/\inserttotalframenumber\hspace*{2ex}
    \end{beamercolorbox}%
  }
}

\setbeamercolor{box1}{bg=red!12,fg=black}

\title{Capital Allocation II}
\subtitle{Estimation Risk, Bet Sizing, and Factor Decomposition}
\author[Moreira]{Alan Moreira}
\institute[UG54]{UG54}
\date{2026}

\begin{document}

% ============================================================
\begin{frame}
\titlepage
\end{frame}

% ============================================================
% SECTION: Capital allocation under uncertainty
% ============================================================

\begin{frame}{Capital allocation under uncertainty}

We know how to maximize the Sharpe ratio given known moments:
$$W^* \propto E[R^e] \, \text{Var}(R^e)^{-1}$$

\vskip 10pt

\textbf{The problem:} in practice, we don't know the true moments.

\vskip 8pt

\begin{itemize}
\item MV optimization loads heavily on assets with extreme sample means and attractive covariance properties
\vskip 6pt
\item We observe only \textit{sample} estimates --- the ex-post tangency portfolio, not the ex-ante one
\vskip 6pt
\item The portfolio that did best in-sample may be a statistical fluke
\end{itemize}

\vskip 10pt

\begin{beamercolorbox}[rounded=true,sep=8pt]{box1}
Whatever your research process, you are likely to end up quite uncertain about the alphas of your trading strategies.
\end{beamercolorbox}

\end{frame}

% ============================================================

\begin{frame}{Capital allocation under uncertainty}

\textbf{Our approach:}

\vskip 8pt

\begin{enumerate}
\item Measure the uncertainty in our risk premia estimates for each asset
\vskip 6pt
\item Show how weights change as we perturb estimates within plausible error bands
\vskip 6pt
\item Show how sensitive the benefits of optimization are to these perturbations
\end{enumerate}

\vskip 10pt

\textit{See notebook: data loading and initial MV weights.}

\end{frame}

% ============================================================
% SECTION: Weights sensitivity to uncertainty
% ============================================================

\begin{frame}{Weights sensitivity to uncertainty}

\textbf{Experiment:} perturb each asset's expected return by $\pm 2$ standard errors of the sample mean, one at a time. Recompute MV weights.

\vskip 10pt

The standard error of the sample mean:
$$\text{SE}(\hat{\mu}_i) = \frac{\hat{\sigma}_i}{\sqrt{T}}$$

\vskip 10pt

\begin{beamercolorbox}[rounded=true,sep=8pt]{box1}
Small changes in inputs can flip signs of portfolio weights.

Large uncertainty means your weights can be way off.
\end{beamercolorbox}

\vskip 10pt

\textit{See notebook: perturbation loop over assets, weight table.}

\end{frame}

% ============================================================
% SECTION: Performance sensitivity to uncertainty
% ============================================================

\begin{frame}{Performance sensitivity to uncertainty}

Compute the Sharpe ratio for each perturbed weight vector (evaluated at sample moments):

\vskip 8pt

$$\text{SR} = \frac{W' \hat{\mu}}{\sqrt{W' \hat{\Sigma} \, W}} \times \sqrt{12}$$

\vskip 10pt

\begin{itemize}
\item This only perturbs \textit{one} mean at a time
\item In practice, you are also off on correlations, variances, and \textit{several} means simultaneously
\item With a long sample and constant-mean assumption, standard errors understate true uncertainty
\end{itemize}

\vskip 10pt

The real test: \textbf{out-of-sample performance} (coming in Performance Evaluation lecture).

\vskip 5pt

\textit{See notebook: SR sensitivity table.}

\end{frame}

% ============================================================
% SECTION: Bet sizing under uncertainty
% ============================================================

\begin{frame}{Bet sizing under uncertainty}

MV optimization is sensitive to inputs $\Rightarrow$ the industry developed many \textbf{ad-hoc bet-sizing rules}.

\vskip 10pt

Bet sizing is one of the great skills in portfolio management --- it requires instinct for uncertainty.

\vskip 10pt

All rules share the structure: compute relative weights $W_i$, then scale by $w$ to hit a volatility target.

\vskip 10pt

\textbf{Best applied after hedging common factors} (market, size, value, etc.) --- i.e., to residual ``alpha'' strategies where $\alpha_i$ is the expected excess return.

\end{frame}

% ============================================================

\begin{frame}{Mean-Variance rule}

$$W_i = w \; \alpha_i \; [\Sigma_\varepsilon^{-1}]_i$$

\vskip 8pt

If strategies are uncorrelated:
$$W_i = w \; \frac{\alpha_i}{\sigma_{\varepsilon,i}^2}$$

\vskip 10pt

\begin{itemize}
\item Optimal when moments are known
\item In practice, highly sensitive to estimation error in $\alpha$
\end{itemize}

\end{frame}

% ============================================================

\begin{frame}{$1/N$ rule}

$$W_i = \frac{1}{N} \; \text{sign}(\alpha_i)$$

\vskip 10pt

\begin{itemize}
\item Ignore the magnitude of alpha --- just bet on the direction
\item Good if you have reliable sign information but noisy magnitude estimates
\item Surprisingly hard to beat out of sample (DeMiguel, Garlappi \& Uppal, 2009)
\end{itemize}

\end{frame}

% ============================================================

\begin{frame}{Proportional rule}

$$W_i = w \; \alpha_i$$

\vskip 10pt

\begin{itemize}
\item Weight proportional to expected return
\item Ignores variance information entirely
\item May be relevant if you cannot take leverage (long-only constraint)
\end{itemize}

\end{frame}

% ============================================================

\begin{frame}{Risk Parity rule}

Assume all strategies have the same appraisal ratio: $\;\dfrac{\alpha_i}{\sigma_i} = \dfrac{\alpha_j}{\sigma_j}$

\vskip 8pt

Substituting into the MV rule:
$$W_i = w \; \frac{1}{\sigma_{\varepsilon,i}}$$

\vskip 8pt

This equalizes the volatility contribution of each strategy:
$$W_i \, \sigma_{\varepsilon,i} = W_j \, \sigma_{\varepsilon,j}$$

\vskip 10pt

\begin{itemize}
\item Robust when you trust variance estimates but not return estimates
\item Very popular in practice (Bridgewater, AQR)
\end{itemize}

\end{frame}

% ============================================================

\begin{frame}{Minimum Variance rule}

$$W_i = w \; [\Sigma_\varepsilon^{-1} \mathbf{1}]_i$$

\vskip 8pt

If uncorrelated: $\;W_i = w / \sigma_{\varepsilon,i}^2$

\vskip 10pt

\begin{itemize}
\item Assumes all alphas are the same --- uses only variance information
\item Solution to: $\;\min_W \text{Var}(W' R^e) \;\;\text{s.t.}\;\; \mathbf{1}'W = 1$
\item Flip strategies so all have positive alpha before applying
\end{itemize}

\end{frame}

% ============================================================

\begin{frame}{Variance Shrinkage rule}

Shrink $\Sigma$ toward a simpler target:
$$\hat{\Sigma}_\tau = (1 - \tau) \, \Sigma + \tau \, \bar{\sigma}^2 \, I$$

\vskip 5pt

where $\tau \in [0,1]$ and $\bar{\sigma}^2$ is the average idiosyncratic variance.

\vskip 8pt

Then apply MV or Min-Variance with $\hat{\Sigma}_\tau$:

$$W_i = w \; \alpha \, \hat{\Sigma}_\tau^{-1}[i] \qquad \text{or} \qquad W_i = w \; \mathbf{1} \, \hat{\Sigma}_\tau^{-1}[i]$$

\vskip 10pt

\begin{itemize}
\item Shrinks all variances toward the average, all correlations toward zero
\item Stabilizes the covariance matrix --- fewer extreme weights
\item Generally a good idea; in some sense this \textit{is} MV with a better $\Sigma$ estimator
\end{itemize}

\vskip 5pt

\textit{See notebook: all six rules applied to FF6 factors, targeted at 16\% annual vol.}

\end{frame}

% ============================================================
% SECTION: Portfolios and Factor models
% ============================================================

\begin{frame}{Portfolios and Factor models}

Given $N$ assets with excess returns $R$, a factor model decomposes:
$$R = A + B f + U$$

\vskip 5pt

where $A$ = alphas, $B$ = loadings, $f$ = factor return, $U$ = residuals ($E[U] = 0$).

\vskip 10pt

A portfolio with weights $W$ has return:
$$r_p = W' R = W'(A + Bf + U)$$

\vskip 10pt

\textbf{Expected return:}
$$E[r_p] = W' A + (W' B) \, E[f]$$

\end{frame}

% ============================================================

\begin{frame}{Variance decomposition}

$$\text{Var}(r_p) = \underbrace{(W' B)^2 \, \text{Var}(f)}_{\text{systematic risk}} + \underbrace{W' \text{Var}(U) \, W}_{\text{idiosyncratic risk}}$$

\vskip 10pt

\begin{itemize}
\item \textbf{Systematic share:} $(W'B)^2 \text{Var}(f) \;/\; \text{Var}(r_p)$
\vskip 6pt
\item Why not just use the in-sample variance of $r_p$ directly?

$\Rightarrow$ Factor structure lets you estimate the covariance matrix with far fewer parameters --- critical when $N$ is large
\vskip 6pt
\item \textbf{Key assumption that fails:} residuals are uncorrelated across assets

$\Rightarrow$ In practice, use multiple factors or allow for residual correlation structure
\end{itemize}

\vskip 5pt

\textit{See notebook: CAPM regressions for SMB, HML, RMW, CMA, MOM; residual covariance matrix.}

\end{frame}

% ============================================================
% SECTION: The rise of "pod shops"
% ============================================================

\begin{frame}{The rise of ``pod shops''}

Traditional hedge funds: principal traders (Soros, Robertson, Tepper, Cohen, \ldots)

\vskip 8pt

Modern structure: \textbf{pod shops} (Citadel, Millennium)
\begin{itemize}
\item Ideas generated at the pod level (5--10 people per pod)
\item Capital allocated by the principal; factor risk hedged centrally
\end{itemize}

\vskip 10pt

\textbf{Why do pod traders work for Citadel instead of running their own fund?}

\end{frame}

% ============================================================

\begin{frame}{Pod shop math --- the $\sqrt{N}$ law}

$N$ pods, each with Sharpe ratio $\text{sr}$ and uncorrelated residuals.

\vskip 8pt

Pool expected return: $\;E\left[\sum W_i^* r_i\right] = N \cdot \text{sr}^2$

\vskip 5pt

Pool variance: $\;\text{Var}\left[\sum W_i^* r_i\right] = N \cdot \text{sr}^2$

\vskip 10pt

\begin{beamercolorbox}[rounded=true,sep=8pt]{box1}
$$\text{SR}_{\text{pool}} = \sqrt{N} \;\cdot\; \text{sr}_{\text{pod}}$$
\end{beamercolorbox}

\vskip 8pt

\begin{itemize}
\item Individually: wealth grows at rate $\sigma \cdot \text{sr}$
\item In the pool: each pod's effective rate is $\sqrt{N} \cdot \sigma \cdot \text{sr}$
\vskip 6pt
\item \textcolor{red}{Finding a new uncorrelated idea is super valuable!}
\vskip 6pt
\item Marginal value of pod $N{+}1$: $\;\sigma \cdot \text{sr}_{\text{pool}} / N$ --- decreasing, but always positive
\end{itemize}

\end{frame}

% ============================================================

\begin{frame}{Pod shops --- caveats}

This calculation \textbf{understates} the value of pooling:
\begin{itemize}
\item Higher SR $\Rightarrow$ more vol capacity without significant loss risk
\item Wealth compounds faster
\end{itemize}

\vskip 10pt

But it is also \textbf{not realistic}:
\begin{itemize}
\item Marginal pods are likely worse (lower SR)
\item Pods are likely correlated with each other (shared factor exposures)
\item Capital constraints and operational costs
\end{itemize}

\end{frame}

% ============================================================
% KEY TAKEAWAYS
% ============================================================

\begin{frame}{Key Takeaways}

\begin{itemize}
\item \textbf{Estimation risk matters.} MV weights are highly sensitive when signals are weak or assets are collinear --- always stress-test inputs.
\vskip 6pt
\item \textbf{Robust bet sizing.} Simple baselines ($1/N$, vol targeting, risk parity) often outperform noisy MV out of sample.
\vskip 6pt
\item \textbf{Shrinkage helps.} Regularizing $\Sigma$ (and sometimes $\mu$) stabilizes weights and improves out-of-sample SR.
\vskip 6pt
\item \textbf{Factor lens.} Decomposing risk by factor clarifies whether diversification is \textit{true} (low common exposures) or \textit{illusory} (correlated bets).
\vskip 6pt
\item \textbf{Pod shops.} Aggregating uncorrelated strategies scales SR by $\sqrt{N}$ --- the key economic logic behind the multi-PM model.
\end{itemize}

\end{frame}

% ============================================================
% APPENDIX
% ============================================================
\appendix
\begin{frame}
\begin{center}
{\Large Appendix}
\end{center}
\end{frame}

\begin{frame}{Bet sizing rules --- summary}

\small
\begin{center}
\begin{tabular}{lll}
\toprule
\textbf{Rule} & \textbf{Formula (uncorrelated)} & \textbf{When to use} \\
\midrule
Mean-Variance & $W_i = w \, \alpha_i / \sigma_i^2$ & Known moments \\
$1/N$ & $W_i = \text{sign}(\alpha_i)/N$ & Reliable sign, noisy magnitude \\
Proportional & $W_i = w \, \alpha_i$ & No leverage, trust means \\
Risk Parity & $W_i = w / \sigma_i$ & Trust vol, not returns \\
Min-Variance & $W_i = w / \sigma_i^2$ & All alphas similar \\
Shrinkage & $W_i = w \, \alpha_i / \hat{\sigma}_{\tau,i}^2$ & Regularize noisy $\Sigma$ \\
\bottomrule
\end{tabular}
\end{center}

\end{frame}

\begin{frame}{References}

\small
\begin{itemize}
\item DeMiguel, Garlappi \& Uppal (2009) --- Optimal versus naive diversification
\item Ledoit \& Wolf (2004) --- Honey, I shrunk the sample covariance matrix
\item Black \& Litterman (1992) --- Global portfolio optimization
\item Asness, Frazzini \& Pedersen (2012) --- Leverage aversion and risk parity
\end{itemize}

\end{frame}

\end{document}
